\section{Задание}

Используя хвостовую рекурсию, разработать программу, позволяющую найти:

\begin{enumerate}
	\item n!,
	\item n-e число Фибоначчи.
\end{enumerate}

Убедиться в правильности результатов. Для одного из вариантов ВОПРОСА и каждого задания составить таблицу, отражающую конкретный порядок работы системы:

Вопрос: \dots

\begin{tabularx}{\textwidth}{|p{1.1cm}|X|X|X|}
	\hline
	№ шага & Состояние резольвенты, и вывод: дальнейшие действия (почему?) & Для каких термов запускается алгоритм унификации: Т1=Т2 и каков результат (и подстановка) & Дальнейшие действия: прямой ход или откат (почему и к чему приводит?) \\
	\hline
	1....  & ...                                                           &                                                                                           & Комментарий, вывод...                                                 \\
	\hline
	2      & ...                                                           &                                                                                           & ...                                                                   \\
	\hline
\end{tabularx}

Т.к. резольвента хранится в виде стека, то состояние резольвенты требуется отображать в столбик: вершина -- сверху! Новый шаг надо начинать с нового состояния резольвенты!

Для одного из вариантов ВОПРОСА и конкретной БЗ составить таблицу, отражающую конкретный порядок работы системы, с объяснениями: очередная проблема на каждом шаге и метод ее решения; каково новое текущее состояние резольвенты, как получено; какие дальнейшие действия? (Запускается ли алгоритм унификации? Каких термов? Почему этих?); вывод по результатам очередного шага и дальнейшие действия.
